\documentclass[onecolumn]{article}

\usepackage{tikz}
\usetikzlibrary{arrows}

\usepackage{graphicx}
\usepackage{verbatim}
\usepackage{subfigure}
\usepackage{multicol}
\usepackage{hyperref}
\usepackage[T1]{fontenc}
\usepackage[utf8] {inputenc} %tildes

\usepackage[all]{xy}
\usepackage{amsmath}
\usepackage{amssymb}
\usepackage{amsfonts}
\usepackage{lscape}
\usepackage{amsthm}
\usepackage[mathscr]{euscript}

\usepackage{fancyhdr}
\usepackage{multirow}
%\usepackage{booktabs}

%\usepackage{tikz}


\usepackage{layout}
\setlength{\hoffset} {0cm}        % 1
\setlength{\voffset} {0cm}        % 2
\setlength{\oddsidemargin}{0cm}   % 3
\setlength{\evensidemargin}{0cm}  % 4
\setlength{\topmargin}{1cm}       % 5
\setlength{\headheight}{4cm}      % 6
\setlength{\headsep}{1cm}         % 7
\setlength{\textwidth}{17cm}      % 8
\setlength{\textheight}{17cm}     % 9
\setlength{\marginparsep}{0cm}    % 10
\setlength{\marginparwidth}{0cm}  % 11
\setlength{\footskip}{0cm}        % 12


%%% Definiciones para su Gu�a de laboratorio%%%
\newcommand{\profesor}{Leonardo Rodríguez Mújica}
\newcommand{\correo}{lrodriguez@mail.unicundi.edu.co}
\newcommand{\nucleotematico}{Antenas y líneas de transmisión}
\newcommand{\codigo}{}
\newcommand{\monitor}{Andrea Carolina Garay Mosquera}
\newcommand{\cormon}{a.cgaray@hotmail.com}

\newcommand{\marco}[2]
{
\begin{center}
{\fboxsep 10pt
\fcolorbox{black}{white}{
\begin{minipage}[t]{16cm}

\section{#1}
#2
\end{minipage}
}}
\end{center}
}

\newcommand{\objetivos}[1] { \marco{Objetivos}{#1} }
\newcommand{\materiales}[1] { \marco{Materiales sugeridos para construcción antena Yagi}{#1} }
\newcommand{\planteamiento}[1] { \marco{Planteamientos y procedimiento}{#1} }
\newcommand{\metodologia}[1] { \marco{Metodología y evaluación}{#1} }
\newcommand{\recomendaciones}[1] { \marco{Recomendaciones}{#1} }
\newcommand{\bibliografia}[1] { \marco{Bibliografia}{#1} }

%%% Digitado por el usuario {\tt }
%%% Click en alguna parte {\sl }
%%% ingl�s {\em }

\cfoot{
Diagonal 18 No 20-29  Fusagasugá- Cundinamarca Teléfonos: 8732512/35 - Fax \\
Página Web: \href{www.unicundi.edu.co}{
www.unicundi.edu.co} - Correo Electrónico \href{mailto:unicundi@unicundi.edu.co}{unicundi@unicundi.edu.co}
}

\fancyhead{
\begin{tabular}{|c|c|c|c|}
\hline
  % after \\: \hline or \cline{col1-col2} \cline{col3-col4} ...
  \hspace{3cm}  & \hspace{8cm}                        & \hspace{2cm}    & \hspace{2cm}     \\
                & \textbf{UNIVERSIDAD DE CUNDINAMARCA}  & FECHA DE          &  CÓDIGO            \\
                & \textbf{FACULTAD DE INGENIERÍA}       & EMISIÓN           &                    \\
                & \textbf{INGENIERÍA ELECTRÓNICA}       &                   &                    \\
                &                                       &  \date{\today}    &                    \\ \cline{2-4}
                &                                       &                   &                    \\
                & \textbf{GUÍA DE LABORATORIO  }        & VERSION No.1.1    &  PÁGINA \thepage   \\
                &                                       &                   &                    \\
  \hline
\end{tabular}
}

\fancyhead[L]{
\hspace{0.4cm}
\vspace{0.4cm}
\includegraphics[scale=0.12]{logo.eps}
}

\title{\textbf{FORMATO GUÍA DE LABORATORIO \\ UNIVERSIDAD DE CUNDINAMARCA \\ FACULTAD DE INGENIERÍA \\ PROGRAMA INGENIERIA ELECTRÓNICA}}
\author{\textbf{Jose David Alvarado Moreno} \\ \\ \huge{\textbf{\LaTeX}}}
\date{\today}

\begin{document}

\pagestyle{fancy}

%%\maketitle

\begin{center}
\begin{tabular}{|l l|l l|}
\hline
\multirow{2}{*}{Asignatura:} &\multirow{2}{*}{\textbf{\nucleotematico}}  & \multirow{2}{*}{Código:}  & \multirow{2}{*}{\codigo}  \\
                             &                                         &                          &                    \\ \hline
\multirow{2}{*}{Docente:}    &\multirow{2}{*}{\textbf{\profesor}}   & \multirow{2}{*}{E-mail:} & \multirow{2}{*}{\href{mailto:\correo}{\correo}}  \\
                             &                                         &                          &                     \\ \hline
\multirow{2}{*}{Monitor:}    &\multirow{2}{*}{\textbf{\monitor}} & \multirow{2}{*}{E-mail:} & \multirow{2}{*}{\href{mailto:\cormon}{\cormon}}   \\
                             &\hspace{6cm}                             &                          & \hspace{6cm}        \\ \hline
\end{tabular}
\end{center}

\begin{center}
\section*{PRACTICA No 3}
\addcontentsline{toc}{section}{PRÁCTICA No 3}

\section*{Implementacion y pruebas antena YAGI}

\addcontentsline{toc}{section}{construccion antena}
\end{center}
%%%--- Colocar objetivos acá ----%%%%

\objetivos{

\begin{itemize}
\item  Comprobar funcionamiento de la antena implementada
\item  Crear una red ad-hoc
\item Realizar y analizar diferentes pruebas antena implementada
\end{itemize}
}

%%%----- Listar los materiales acá ------%%%
\materiales{

\begin{enumerate}
  \item dos PC
  \item Antena Yagi-Uda
  \item Pig tail (cola de cerdo)
  \item Tarjeta de red (150n)
\end{enumerate}
}


%%% ------- colocar planteamiento y procedimiento

%% 1. simular la yagi con yagi model
%% 2. con las fotos mostrar como debe quedar
%% 3. comparar con la simulación de xnec2c


%% 1.  hacer el balum
%% averiguar bien el tipo de cable con que se va a hacer
%% 2. probar la antena con el equipo
%% 3. enseñar a conectarse en linux
%% 4. ver la intensidad de la señal con wavemon


%% grupo especial para práctica de patron de radiación. Replicar la práctica de méxico.
%% 1. probar dipolo en equipo
%% 2. probar galvanometro con diferentes access points y equipos de forma incremental.

\planteamiento{
\begin{enumerate}

%% red ad-hoc
\item  Conéctese a una red inalámbrica con el mono polo de la tarjeta de red ¿en qué modo está trabajando la tarjeta de red?

%% red ad-hoc
\item Cree una red ad-hoc entre dos equipos Linux con los siguientes pasos:

\begin{enumerate}

    \item Tal vez tenga que deshabilitar el administrador de la red en Linux e Identifique la interfaz inalámbrica con que va a trabajar con las órdenes 
    
    {\tt \# /etc/init.d/network-manager stop}
    
    {\tt \# iwconfig}
    
\item Deshabilite la interfaz de red. Según el dispositivo que utilicemos habrá que especificar la entrada correcta.

{\tt \# ifconfig wlan0 down}

\item Pase a modo ad-hoc la tarjeta de red. 

{\tt \# iwconfig  wlan0 mode ad-hoc}

\item Le asignamos un nombre a nuestra red

{\tt \# iwconfig wlan0 essid "Mired"}

%\item Luego de ello procedemos a configurar el canal y es por medio del cual fluirán los datos. Si existen más redes en el área les recomendaría utilizar un canal que este libre.

%{\tt \# iwconfig wlan0 channel x}

\item Asignamos una dirección IP a la máquina. :

{\tt \# ifconfig wlan0 192.162.0.1 netmask 255.255.255.0}

\item En la otra máquina debe examinar si la red se está viendo,  conectese a la red y asigne una IP que pertenezca a la misma red

{\tt \# iwlist wlan0 scan }

{\tt \#  	iwconfig wlan0 essid  "Mired"}

\item Ahora procedemos a habilitar el redireccionamiento. Ejecutamos el siguiente comando:

{\tt \#echo 1>proc /sys/net/ipv4/ip_forward}

\item pruebe la comunicación con {\tt ping}

\item Si no funciona luego del tercer intento, haga la comunicación a través de {\tt netwotk-manager} (sobre el ícono --> botón secundario del ratón --> editar conexiones)
\end{enumerate}

\end{enumerate}

}

\planteamiento{
\begin{enumerate}

%% pruebas
\item Instale la wavemon ({\tt apt-get install wavemon} para ver intensidad de la señal.

\item Pruebe la antena con el equipo

\item Mida la intensidad de la señal a una distancia de referencia. 
Tome 10 puntos de la señal alejándose de la fuente, a distancias aproximadamente uniforme. Examine la intensidad de la señal con la herramienta wavemon

\item Realice un gráfico aproximado de distancia metros (m) vs potencia (W).

\item Con la potencia a la distancia de referecia, calcule las potencias de los otros 10 puntos como si esta potencia fuera la potencia del transmisor




\end{enumerate}
}

%%% --- Colocar acá la metodología y evaluación
\metodologia{

Presentar un informe con los resultados obtenidos en la práctica que se
evaluará  teniendo en cuenta la siguiente rúbrica.

\newcommand{\med}{0.22}
\begin{tabular}{|p{\med\textwidth}|p{\med\textwidth}|p{\med\textwidth}|p{\med\textwidth}|p{\med\textwidth}|}
\hline
criterio&Excelente(90\%)&Bueno(70\%)&Deficiente(10\%)\\
\hline
Pruebas con la antena Yagi(25\%)&Documentación de las pruebasna&Genera documento sin las principales tareas&No documenta pruebas de la antena\\
\hline
Funcionamiento de la antena Yagi(25\%)&Realiza y analiza pruebas con la antena&Implementa pero no realiza pruebas&No obtienes resultados necesarios del funcionamiento de la antena\\
\hline
Crea red ADHOC(25\%)&Crea y realiza pruebas con la red ADHoC&realiza pruebas deficientes&No realiza las actividades\\
\hline
Forma:presentación, redacción ortografía (10\%)&Sin errores y bien
presentado&Presenta deficiencias&No elabora documento\\
\hline
Conclusiones(15\%)&claras, concisas y relevantes&Falta especificar resultados&Irrelevantes e intrasendentes\\
\hline
\end{tabular}
}
%%% --- colocar acá las recomendaciones ---%%%
\recomendaciones{

\begin{itemize}
\item Leer e investigar sobre como crear red ADHOC


\end{itemize}
}

%%% ---- listar las referencia bibliogr�ficas
\bibliografia{

\begin{itemize}
  \item Xnec2c User Manual.\url{http://www.qsl.net/5b4az/pkg/nec2/xnec2c/doc/xnec2c.html }
  \item Antena Yagi
Networks.\url{http://arieldx.tripod.com/estaciondx/proyectos/yagifm.htm}
\item Software YagiModel.\url {http://fermi.la.asu.edu/w9cf/yagipub/}
 %% \item \url{http://tf.nist.gov/general/pdf/451.pdf} \url{http://www.researchgate.net/post/Regarding_Yagi-Uda_antenna}
 
\end{itemize}
}




\end{document} 
